Una sostanza pura è una sostanza in cui la composizione chimica non
varia in tutta la massa considerata. \#\#\# Fasi di una sostanza pura
Una sostanza pura si può trovare in 3 fasi: - Solida - Liquida - Gassosa

Un liquido può essere \textbf{sottoraffreddato} (non in procinto di
evaporare) o \textbf{saturo} (in procinto di evaporare). Analogamente un
vapore \textbf{saturo} è in procinto di condensare e uno
\textbf{surriscaldato} no. Tra un liquido saturo e un vapore saturo c'è
la \textbf{miscela satura di liquido e vapore}.
\includegraphics{Fasi.png}

\emph{Diagramma di un processo di riscaldamento dell'acqua a pressione
costante}

\hypertarget{proprietuxe0}{%
\subsubsection{Proprietà}\label{proprietuxe0}}

\textbf{Temperatura di saturazione} \(T_\text{sat}\): Temperatura a cui
una sostanza inizia a evaporare (fissata la pressione) \textbf{Pressione
di saturazione} \(p_\text{sat}\): Pressione a cui una sostanza inizia a
evaporare (fissata la temperatura) \textbf{Calore latente}: quantità di
energia assorbita (o liberata) durante trasformazione con cambio di fase
\textbf{Calore latente di fusione}: La quantità di energia assorbita
durante la fusione uguale a quella liberata durante solidificazione.
\textbf{Calore latente di vaporizzazione}: La quantità di energia
assorbita durante la vaporizzazione uguale a quella liberata durante la
condensazione. \textbf{Punto critico}: punto in cui coesistono vapore
saturo e liquido saturo. \textbf{Volume specifico} \(v\): reciproco
della densità \textbf{Titolo} \(x\): rapporto tra massa vapore e massa
totale in una miscela. \(x\in[0,1]\), se \(x=0\) liquido saturo, se
\(x=1\) vapore saturo. Un liquido saturo mantiene le proprietà anche
quando è in una miscela con vapore.

\hypertarget{entalpia}{%
\subsubsection{Entalpia}\label{entalpia}}

\(h=u+pv\hspace{0.4cm}(kJ/kg)\) \(H=U+pV\hspace{0.4cm}(kJ)\) Entalpia di
vaporizzazione \(h_\text{l,sat}\) (calore latente di vaporizzazione):
Energia necessaria a trasformare l'unità di massa di un liquido saturo
in vapore saturo a una certa pressione e temperatura.
